\section*{Bab \ref{ch:g9:inside-the-atom} --- Di Dalam Atom}

\begin{solution}{exo:g9:inside-the-atom:1}
Di dalam inti: \omterm{def:g9:inside-the-atom:proton}{proton} (positif) dan \omterm{def:g9:inside-the-atom:proton}{neutron} (tidak bermuatan). Di
sekelilingnya: \omterm{def:g9:inside-the-atom:nucleus}{elektron} (negatif).
\end{solution}

\begin{solution}{exo:g9:inside-the-atom:2}
\ce{^{16}_{8}O}: 8 \omterm{def:g9:inside-the-atom:proton}{proton}, 8 \omterm{def:g9:inside-the-atom:proton}{neutron}, 8 \omterm{def:g9:inside-the-atom:nucleus}{elektron}.
\ce{^{27}_{13}Al}: 13, 14, 13. \ce{^{40}_{20}Ca}: 20, 20, 20.
\end{solution}

\begin{solution}{exo:g9:inside-the-atom:3}
\omterm{def:g7:atoms-and-molecules:atom}{Atom} itu memiliki \omterm{def:g9:inside-the-atom:nucleus}{elektron} (masing-masing bermuatan $-e$) sebanyak
\omterm{def:g9:inside-the-atom:proton}{protonnya} (masing-masing bermuatan $+e$): muatan-muatannya saling
meniadakan.
\end{solution}

\begin{solution}{exo:g9:inside-the-atom:4}
$A = 26 + 30 = 56$: \ce{^{56}_{26}Fe}.
\end{solution}

\begin{solution}{exo:g9:inside-the-atom:5}
Ya: keduanya memiliki $Z = 17$, jadi keduanya klorin. \omterm{def:g9:inside-the-atom:atomic-number}{Nomor massanya}
berbeda (18 dan 20 \omterm{def:g9:inside-the-atom:proton}{neutron}): keduanya \omterm{def:g9:inside-the-atom:isotope}{isotop}.
\end{solution}

\begin{solution}{exo:g9:inside-the-atom:6}
Bukan: \omterm{def:g9:inside-the-atom:atomic-number}{nomor atomnya} berbeda (6 dan 7), jadi keduanya unsur yang
berbeda, karbon dan nitrogen. \omterm{def:g9:inside-the-atom:isotope}{Isotop} memiliki $Z$ yang sama.
\end{solution}

\begin{solution}{exo:g9:inside-the-atom:7}
$56 \times 1.67 \times 10^{-27} \approx \qty{9.35e-26}{kg}$.
\end{solution}

\begin{solution}{exo:g9:inside-the-atom:8}
\omterm{def:g7:atoms-and-molecules:atom}{Atom} sebagian besar berupa ruang kosong: intinya sangat kecil, sehingga
kebanyakan partikel tidak bertemu inti mana pun dalam perjalanannya dan
lewat lurus.
\end{solution}

\begin{solution}{exo:g9:inside-the-atom:9}
$10^{-10} / 10^{-15} = 10^{5}$: seratus ribu kali lebih kecil.
\end{solution}

\begin{solution}{exo:g9:inside-the-atom:10}
Sifat kimia bergantung pada \omterm{def:g9:inside-the-atom:nucleus}{elektron}, dan \omterm{def:g9:inside-the-atom:isotope}{isotop-isotop} sebuah unsur
memiliki jumlah \omterm{def:g9:inside-the-atom:nucleus}{elektron} yang sama ($Z$ yang sama).
\end{solution}

\begin{solution}{exo:g9:inside-the-atom:11}
$26 \times 9.11 \times 10^{-31} = \qty{2.37e-29}{kg}$. Bagiannya:
$2.37 \times 10^{-29} / 9.35 \times 10^{-26} \approx 2.5 \times
10^{-4}$, sekitar \qty{0.025}{\percent} massanya.
\end{solution}

\begin{solution}{exo:g9:inside-the-atom:12}
6915 \omterm{def:g7:atoms-and-molecules:atom}{atom} tembaga-63 dan 3085 \omterm{def:g7:atoms-and-molecules:atom}{atom} tembaga-65. Rata-rata:
$(6915 \times 63 + 3085 \times 65)/10\,000 = 63.6$ \omterm{def:g9:inside-the-atom:proton}{nukleon}.
\end{solution}

\begin{solution}{pb:g9:inside-the-atom:1}
\textbf{1.} \ce{^{35}_{17}Cl} dan \ce{^{37}_{17}Cl}.

\textbf{2.} Klorin-35: 17 \omterm{def:g9:inside-the-atom:proton}{proton}, 18 \omterm{def:g9:inside-the-atom:proton}{neutron}, 17 \omterm{def:g9:inside-the-atom:nucleus}{elektron}.
Klorin-37: 17 \omterm{def:g9:inside-the-atom:proton}{proton}, 20 \omterm{def:g9:inside-the-atom:proton}{neutron}, 17 \omterm{def:g9:inside-the-atom:nucleus}{elektron}.

\textbf{3.} Keduanya memiliki 17 \omterm{def:g9:inside-the-atom:proton}{proton}, jadi keduanya unsur klorin;
keduanya hanya berbeda dalam banyaknya \omterm{def:g9:inside-the-atom:proton}{neutron}, jadi keduanya \omterm{def:g9:inside-the-atom:isotope}{isotop}.

\textbf{4.} Ya: sifat kimia bergantung pada \omterm{def:g9:inside-the-atom:nucleus}{elektron}, dan keduanya
memiliki 17 \omterm{def:g9:inside-the-atom:nucleus}{elektron}.

\textbf{5.} $35 \times 1.67 \times 10^{-27} = \qty{5.845e-26}{kg}$.

\textbf{6.} $17 \times 9.11 \times 10^{-31} = \qty{1.55e-29}{kg}$:
sekitar 4000 kali lebih kecil daripada intinya, dapat diabaikan.

\textbf{7.} $37 \times 1.67 \times 10^{-27} = \qty{6.179e-26}{kg}$.

\textbf{8.} $\qty{1}{g} = \qty{e-3}{kg}$; $10^{-3} / 5.845 \times
10^{-26} \approx 1.71 \times 10^{22}$ \omterm{def:g7:atoms-and-molecules:atom}{atom}.

\textbf{9.} \qty{75.8}{\percent} dan \qty{24.2}{\percent}.

\textbf{10.} $758 \times 5.845 \times 10^{-26} = \qty{4.430e-23}{kg}$;
$242 \times 6.179 \times 10^{-26} = \qty{1.495e-23}{kg}$.

\textbf{11.} $4.430 \times 10^{-23} + 1.495 \times 10^{-23} =
\qty{5.925e-23}{kg}$.

\textbf{12.} $5.925 \times 10^{-23} / 1000 \approx
\qty{5.93e-26}{kg}$: massa rata-rata sebuah \omterm{def:g7:atoms-and-molecules:atom}{atom} klorin.
\end{solution}
